% ---------------------------------------------------------------------------
% Related work
%
% Frontier-Eng (arXiv 2604.12290) goes first and gets the most room: its
% formalization is nearly ours. A task is (C, x0, E) -- read-only context, a
% feasible-but-naive starting artifact, a frozen evaluator returning a
% feasibility flag plus a scalar score -- and the agent proposes revisions from
% the full history under a fixed budget. Three differences to state plainly,
% because they are what this paper adds rather than repeats:
%   1. it has no human baseline and no RL baseline; its reference points are
%      contributor-written naive programs plus correctness anchors. Work items
%      2 and 3 are exactly those two missing cells.
%   2. its budget is evaluator interactions only -- no tokens, dollars, or
%      wall-clock.
%   3. it optimizes against a frozen evaluator with no opponent, so its score is
%      absolute and monotone. Elo is neither, which is why its metrics cannot be
%      adopted directly: a rating is relative to a population, and a gain can come
%      from the field weakening. Be precise about which pool is meant -- HL's test
%      pool is frozen (that is what makes our Elo comparable within a run), while
%      the *archive's* human population drifts across editions (which is what the
%      cross-edition anchor exists to absorb). Do not write "drifting opponent
%      pool" without saying which, and never of the frozen test pool.
% Cite its numbers as reported: the abstract says eight models and the
% experiments nine, the power law rests on one model-framework pair, and the
% depth/width result on a 10-task subset.
%
% MLE-bench (arXiv 2410.07095) for the premise that agents iterate programs
% effectively; arXiv 2507.02554 for the cleaner framing of a research agent as a
% search policy over a solution space, which is how the HL loop should be
% described. Frontier-Eng's depth-vs-width finding -- one deep chain beats
% parallel restarts at fixed budget, because restarts discard accumulated
% context -- is the direct evidence for the long-trajectory premise.
%
% Programmatic RL is the nearest representational neighbour and needs its own
% paragraph: programmaticrl2024foundations for the size/expressiveness tradeoff
% (it bounds the size of optimal programmatic policies, which is the formal
% relative of K(pi_good) -- cite it for that and not for editability, which it
% never claims), diprl2026 and llmguidedsearch2024 for the current line,
% kuang2025generative for
% the closest method (LLM revises a Python policy from execution traces, beats
% deep RL baselines at far fewer interactions). What we add over kuang2025generative:
% adversarial rather than single-agent, a human field to be placed within, a
% shared budget ledger, and a ceiling analysis. Their efficiency claim is
% directional with no normalization protocol, which is exactly the gap our
% ledger fills -- say that without disparaging the work.
%
% SlopCodeBench (slopcodebench2026) belongs here as well as in the introduction,
% and the paragraph has to draw the distinction the introduction draws: its
% degradation is driven by a *growing specification*, so accumulation is
% compulsory. A game's rules are fixed, so a better heuristic can replace a worse
% one and the program may shrink. Without that distinction a referee reads their
% result as predicting our loop cannot work.
%
% agent2rlbench2026 earns one sentence, for a point that is easy to miss: in
% AI4AI the agent is already a symbolic program manipulating a connectionist
% artifact, so the first thing an LLM does to RL is write code around it. That is
% what makes "then why not have it write the policy itself" the obvious next
% question rather than a contrarian one -- which is this paper. Do not overstate
% it into a claim that they anticipated HL; they benchmark agents engineering RL
% post-training, a different task.
%
% Then: RL in games, program synthesis / LLM-driven code search, automated
% heuristic design, continual learning. The 21-item survey in the earlier
% proposal is the starting list (that PDF is not in this repo -- it was never
% tracked, so treat refs.bib as the authoritative list and re-obtain the
% proposal from the user if the original 21 items are needed).
%
% SURVEYED AND DELIBERATELY NOT CITED. Recorded so a later pass does not spend
% the search again, and so the rejections are reviewable rather than silent.
%   - arXiv 2604.21896, "Crafting Strategic AI Gaming Agents...": operationalizes
%     Shannon's classification into an authoring environment, with a
%     heuristic-games class built on minimax plus crowd-sourced data. Sounds
%     adjacent and is not: no human baseline, no RL baseline, no step counts, no
%     ceiling. Systems-and-pedagogy contribution; nothing to compare against.
%   - arXiv 2602.04254, "Scaling Agentic Verifier for Competitive Coding": no
%     games and no adversary that adapts. Worth one line ONLY if a referee
%     argues iterative LLM improvement does not plateau, since it reports clear
%     test-time scaling with no saturation. The boundary to state if so: their
%     scaling is over K candidate solutions at test time with a verifier that
%     grows, whereas our claim is about maintaining ONE artifact across
%     revisions. Different quantity scaled, so it is not a counterexample --
%     but do not pretend the tension is absent either.
% Two hits from that same search were kept: echo2025hindsight and tig2025think,
% both cited in Section 1 and both load-bearing there. See their refs.bib
% annotations before moving either.
%
% Note the name clash: AgentBench (arXiv 2308.03688) is a different paper, also
% from Tsinghua, and is the style reference in _reference_paper/. The project
% name is still undecided.
%
% COMPRESSED TO THREE PARAGRAPHS on 2026-08-19 at the user's direction ("related work
% 压缩成3个自然段，并且scale up引用的文章的数量"). Five paragraphs became three, and the
% citation count went from 35 to 45 keys -- every key in refs.bib is now cited somewhere
% in the paper, so a further scale-up requires new bibliography entries and must not be
% done by re-citing existing ones. The three-way split is:
%   1. programs as policies      -- representation, heuristic design, program search,
%                                   and code-written-around-a-learner, merged because
%                                   they share one axis: who writes the program.
%   2. adversarial evaluation    -- games as testbed, pool design, human fields as
%                                   percentile references, and the baseline rationale.
%   3. continual improvement     -- experience reuse AND the whole maintenance-capacity
%                                   argument, formerly three paragraphs of its own.
%
% THE UPPER-BOUND ARGUMENT MOVED HERE on 2026-08-18 (second cut, to reach 2.5 pages) and
% the header note that stood here said: do not compress it into a literature-survey
% paragraph, because the mechanism has to be argued rather than cited or the intro's
% central claim is unsupported anywhere in the paper. That note was OVERRIDDEN by the
% compression above, and the constraint behind it was preserved rather than dropped: the
% mechanism is still STATED IN FULL before its first citation, in paragraph 3, and all
% six load-bearing elements below survive. What was cut is the numbers, not the argument
% -- per-repo figures, the four per-model execution horizons, the exact prompting deltas,
% and the erosion/verbosity slope pairs. They remain recorded in this header, which is
% now their only home in the repo, so verify against it rather than re-reading the
% papers. If the argument ever needs to grow back, grow it in Section 6, not here.
% What paragraph 3 carries, and must keep carrying:
%   - the mechanism itself, stated before any citation.
%   - the RATE-not-exemption form of the human comparison. A human exemption would
%     make capacity look like an agent-specific defect instead of a general limit
%     agents reach sooner.
%   - the prompting result, which answers "the plateau is just a prompting defect".
%   - the replacement-vs-accumulation boundary, without which these results read as
%     predicting our loop cannot work.
%   - the reasoning result's direction, which helps us rather than hurting.
%   - selfimprovement2026survey, whose only citation in the paper is here.
% All figures below were verified against full text; refs.bib stays authoritative.
% Cross-check them against the body if either changes.
%   - slopcodebench2026's PROMPTING RESULT is in the body as of 2026-08-18 and
%     answers a referee who says the plateau is a prompting defect rather than a
%     capacity limit. Never write "statistically unchanged" of it: their paper runs
%     no significance test on that experiment, and "slope"/"intercept" appear only
%     in their Appendix D, about the human panel. The body says the intercept moves
%     while the growth rate is left intact, which is what is supported.
%   - slopcodebench2026 (v2: 36 problems, 196 checkpoints, 15 agents): erosion
%     rises in 77% of trajectories, verbosity in 75.5%; cost per checkpoint 2.2x
%     first-to-last "without corresponding improvements to correctness"; agent
%     code 2.3x more verbose and 2.0x more eroded than 473 open-source Python
%     repos; per-revision accumulation 6.6x (verbosity) and 5.0x (erosion)
%     faster. The human panel: 68% of repos end more verbose (median +36%),
%     erosion rises in 53%, slopes 0.0022/0.0053 vs the agents' 0.0144/0.0264 --
%     a RATE gap, never an exemption. Quality-aware prompting: verbosity -27.5
%     to -35.6%, erosion -34.3 to -57.6% at the intercept, growth rate unchanged,
%     price -2.3 pp strict correctness and +12.1% cost. There is NO significance
%     test on that experiment; never write "statistically unchanged".
%   - longcontextbugfix2026: successful agentic trajectories stay under 20-30k
%     tokens. The 64k single-shot result (7% and 0%) is a DIFFERENT model set and
%     not the agentic harness -- report it as their framing does, decomposition
%     vs retrieval, not "padding the context makes it worse".
%   - longhorizonexecution2026: single-turn execution horizons, GPT-5 2176 steps,
%     Claude-4 Sonnet 432, Grok 4 384, Gemini 2.5 Pro 120. Scaling model size
%     INCREASES self-conditioning; thinking removes it entirely (their Result 5).
%     Their thesis argues against diminishing returns -- we cite the mechanism,
%     not the thesis.
%   - frontiereng2026: t^-1 frequency and k^-1 magnitude, both R^2 > 0.83,
%     marginal returns near zero after ~50-100 iterations, tail to 500, 47 tasks.
%     DELIBERATELY NOT IN THE BODY. The user cut this on 2026-08-18: the paragraph
%     now says only that frequency and magnitude both decay, with no exponents, no
%     R^2, no iteration counts, and no task count. Do not "restore" them as an
%     omission -- it was a decision. The figures stay recorded here because the
%     paragraph above still claims a decay shape and a reviewer may ask for it, and
%     because 05-experiments.tex compares our own \val{decay} against theirs; if that
%     comparison is written, the exponent enters the paper THERE, not here.
% ---------------------------------------------------------------------------
\section{Related Work}

\paragraph{Reinforcement learning on adversarial games.}
Symbolic construction won adversarial games first, by search over a hand-authored
evaluation \citep{deepblue2002}. Learned value estimation then displaced it, and did so by
removing one hand-authored component after another --- the evaluation function, then the
human game record, then the game-specific method, then the environment model itself
\citep{dqn2015,alphago2016,alphagozero2017,alphazero2018,muzero2020,efficientzero2021}.
From there the line extended into exactly the settings our suite occupies: imperfect
information, whether by equilibrium computation or by model-free multiagent learning,
multiplayer and non-zero-sum payoffs, hidden hands over long horizons, real-time control
against champion humans, and action spaces that defeat enumeration
\citep{cfr2007,deepstack2017,libratus2018,pluribus2019,deepnash2022,suphx2020,douzero2021,
gtsophy2022,alphastar2019,openaifive2019}. Where a language model has entered these games
it has been a negotiator beside a planner rather than the policy's author
\citep{cicero2022}.
Two results from this line bear on how we evaluate rather than on what we build. Opponent
pools are a design variable and not a detail, and a policy of superhuman average strength
can still carry holes an adversary finds
\citep{psro2017,adversarialgo2023,asro2026coevolution,adversarialheuristics2026} --- which
is why our test pool is frozen and drawn from a whole human field rather than reduced to
one strong reference, and why cost is reported beside strength
\citep{katago2019,mlebench2024,olympiad2025percentile}. This paradigm's standing is not in
question here. What it leaves open is the small-budget regime this paper measures, and
against which learner: comparing a maintained program against policy gradient alone would
confound symbolic against neural with on-policy against off-policy, so the primary baseline
is off-policy and model-based \citep{offpolicyvaluellm2026,samplereuse2022}.

\paragraph{Policies written as programs.}
A policy can be an executable program rather than an opaque network, and before language
models the program was obtained by asking a learner for it --- a sketch filled by policy
search, a decision tree distilled from a trained network so that it can be verified, a
latent space of programs made searchable, symbolic controllers recovered directly
\citep{pirl2018,viper2018,inductivegen2020,leaps2021,symbolicpolicies2021}. The
representation carries a size--expressiveness trade-off that has since been made formal,
bounding the size of an optimal programmatic policy, which is the relative of the
description length this paper measures \citep{programmaticrl2024foundations}. Learning has
also been aimed at code as the artifact rather than as the medium, discovering algorithms
that were then adopted \citep{alphatensor2022,alphadev2023}. What changed with language
models is who writes the program: they synthesize policies directly, search structured
domains, and supply proposals that evolutionary selection, uncertainty balancing, diversity
preservation, or learned teachers then filter
\citep{llmguidedsearch2024,qube2024,hsevo2024,diprl2026,prorl2026,teacherawareheuristic2026,
heuristicsearch2026,latentheuristicsearch2026,automodsat2025}. A parallel branch writes code
around a learner rather than as the learner, from reward functions to the post-training
machinery itself \citep{eureka2023,rewarddesignagent2026,agent2rlbench2026} --- which is
what makes ``then let it write the policy'' the next question rather than a contrarian one,
and the closest game-specific systems take it up with model-controlled program search
\citep{codetoplay2024,kuang2025generative}. In all of these the program is a search output:
extracted, sampled, or selected from a population. In ours it is a maintained artifact,
revised in place across a long sequence, which is what makes its description length a
trajectory rather than a hyperparameter.

\paragraph{Self-evolving agents, and what stops them.}
An agent can improve with no weight update at all, and this literature divides by what the
improvement is written into. Into natural language, as rewritten trajectories, retained
verbal lessons, distilled deployment rules, or refinement from self-generated feedback
\citep{selfrefine2023,reflexion2023,expel2024,echo2025hindsight,learningonthejob2026}; into
the prompt or a text-valued gradient over it \citep{promptbreeder2023,textgrad2024}; into
the scaffolding, by search over agentic designs \citep{adas2024}; and into code, the branch
this paper belongs to, where model-generated variants filtered by an external evaluator gave
way to autonomous proposal and repair, to coding agents that rewrite themselves, to
libraries of executable skills, and to pipelines running a whole research loop
\citep{funsearch2024,stop2023,generativeagents2023,voyager2023,alphaevolve2025,
selfimprovingcodingagent2025,darwingodelmachine2025,aiscientist2024,avo2026,
mendelgoedel2026}. Alongside these, benchmarks now argue that such learning must be
measured by transfer and interference rather than long-context retrieval, and that a
reported gain must be separated from the harness that produced it
\citep{agentcl2026,clbench2026,harnessengineering2026} --- which is why ours is held fixed.
What distinguishes our case within the code branch is one artifact instead of a population:
experience is not retrieved but written into the program, which is what creates the failure
mode this paper measures. The mechanism is accumulation: revisions leave the program larger and more
interconnected \citep{mccabe1976complexity}, and each further revision must place an edit
without disturbing what already works, until past some size placement fails more often
than it succeeds --- not because the paradigm ran out of things to express, but because the
artifact has outgrown what its maintainer can hold, attribute, and change at once. Each
term of that account is now measured outside our setting. Across 196 checkpoints of
iterative tasks, erosion rises in 77\% of trajectories and verbosity in 75.5\%, at 2.2$\times$
the cost per checkpoint and without matching gains in correctness; human repositories are
not exempt but degrade an order of magnitude more slowly, so capacity is a rate rather
than an agent-specific defect \citep{slopcodebench2026}. Quality-aware prompting moves the
intercept and leaves the growth rate intact, at a cost in correctness, which is evidence
for a capacity rather than a prompting artifact. Two mechanisms underneath are not about
code quality: agentic repair succeeds while the context each step accumulates stays small
\citep{longcontextbugfix2026}, and per-step accuracy falls as steps accumulate even when
the plan is given, though reasoning removes the self-conditioning behind it
\citep{longhorizonexecution2026} --- so a plateau reached by a reasoning agent cannot be
attributed to that effect. A program under revision is the case that cannot be decomposed
away, since the artifact \emph{is} the accumulated context. One boundary keeps these
results from predicting that our loop cannot work at all: the degradation they measure is
driven by a specification that keeps growing, so nothing may be discarded, whereas a
game's rules are fixed and a better heuristic may \emph{replace} a worse one. Accumulation
is a tendency here rather than an obligation, which is why the ceiling is worth measuring
instead of assuming. Frequency and magnitude of improvement do both decay on real
engineering tasks \citep{frontiereng2026}, but those tasks are single-agent and carry
neither a human field nor a learned baseline, leaving the position of the symbolic paradigm
against a learned one open; a survey of this literature organizes it by the strength of the
verification signal the loop closes on \citep{selfimprovement2026survey}, at whose strong
end a game engine with a frozen pool sits. We carry the same longitudinal question into
adversarial games and ask whether the plateau is explained by the maintainable description
length of the policy rather than by the complexity of the game's rules.
